% RESEARCH_DRAFT_SHA256: 60510ae75315c2012229390362ca712847fb27b169bbd50d374f0add922f9f1b
% Layout and source-reference wrappers expanded for the web renderer.
\documentclass[11pt]{article}
\usepackage[margin=25mm]{geometry}
\usepackage{fontspec,unicode-math}
\setmainfont{Latin Modern Roman}
\setmathfont{Latin Modern Math}
\usepackage{amsmath,amssymb,mathtools,xurl,hyperref}
\hypersetup{colorlinks=true,urlcolor=blue,linkcolor=blue}
\setcounter{secnumdepth}{0}
\setlength{\parindent}{0pt}
\setlength{\parskip}{5pt}
\setlength{\emergencystretch}{3em}
\newcommand{\sref}[2]{\hyperlink{source:#1:#2}{[S#2]}}
\newcommand{\cataloguescope}[1]{\paragraph{Source record.}#1}
\newcommand{\floor}[1]{\left\lfloor#1\right\rfloor}
\begin{document}
\section*{Research-attempt provenance}
This standalone research attempt retains the complete original definition body below,
including its historical source and preparation notes. Those notes describe the
earlier catalogue preparation; the new attempt and its verification are identified
separately after that body. No original ID, catalogue status, or source statement is changed.

\section{State Recovery from Quantized Nonlinear Trajectories}\label{statement}

\subsection{Definitions and mathematical statement}

\paragraph{Shared notation.}Standard number systems are $\mathbb N,\mathbb Z,\mathbb Q,\mathbb R,\mathbb C$. Specialized notation and index ranges follow the cited source and the definitions below.


$f(S)=\{f(x):x\in S\}$ and $f^{-1}(S)=\{x:f(x)\in S\}$. A quantization region records the set in which a state lies. The recursive sets collect states consistent with forward or backward observation constraints.
\paragraph{Record scope and qualifications.}
The catalogue title is retained as metadata. Named certificates, contracts, bounds or frameworks in that title are not asserted as results of the original source extract.
The supplied registry flags this extraction as unrecoverable. It is displayed as a literal source fragment; the definitions provide context rather than a replacement formal statement.

Source-record limitation: Source flags the statement as unrecoverable: The official report confirms that this is the broad introductory formulation of the single observability problem stated fully in OWR-3394-018. No separate clean statement is promoted; the original record remains publicly visible for provenance.

Related catalogue records: \href{https://mehmetmars7.github.io/Erdosproblems-llm-hunter/problem.html?type=open_problems&id=30001202}{30001202}. Separate source IDs are preserved for their provenance.
The source text below is rendered literally where necessary; damaged formulas and composite extracts are not silently repaired into a new mathematical claim.
The supplied catalogue marks this record as unrecoverable. Its extract and terminology are documented, but a complete original problem has not been reconstructed.
\paragraph{Statement recorded in the supplied source.}
Consider a known nonlinear discrete-time system x(t+1)=f(x(t)) on a compact X\textbackslash{}subset\textbackslash{}mathbb\{R\}\textasciicircum{}n with a finite quantization into regions P\_i. Given only the regions P(0),...,P(T) containing an unknown trajectory, under what conditions on f can one recover nontrivial information about some x(t) beyond x(t)\textbackslash{}in P(t), equivalently make the feasible sets obtained from the quantized observations small? \href{https://doi.org/10.4171/owr/2009/11}{[S2]}
\paragraph{Source status.}
The supplied catalogue labels this source record partially solved; that metadata was checked against the supplied snapshot on 2026-10-11. The label is preserved for provenance and is not an independent verification of present open status or mathematical correctness. The catalogue literature assessment is dated 2026-08-21.
\subsection{Short English statement}
The supplied fragment concerns state Recovery from Quantized Nonlinear Trajectories. Its complete mathematical target is not recovered.
\subsection{Sources}
\begin{itemize}
\item[S1] Ulam AI and the UnsolvedMath Contributors, UnsolvedMath: A Curated Collection of Open Mathematics Problems, catalogue entry. Catalogue curation and metadata: CC BY 4.0. The source extract is retained; definitions, English summaries and qualifications are editorial additions.
\url{https://huggingface.co/datasets/ulamai/UnsolvedMath}.
\item[S2] Control Theory: On the Way to New Application Fields (2009), Oberwolfach Reports, DOI 10.4171/owr/2009/11, pp. 593--680
\url{https://doi.org/10.4171/owr/2009/11}.
\item[S3] Liberzon--Hespanha 2008
\url{https://doi.org/10.3182/20080706-5-KR-1001.01357}.
\item[S4] nonlinear quantized state estimation
\url{https://doi.org/10.1016/j.neucom.2013.03.022}.
\item Source-specific Oberwolfach Reports licences, including CC BY-SA 4.0 where applicable, are retained; see THIRD\_PARTY\_NOTICES.md.
\end{itemize}
\paragraph{Preparation and review.}
The original source and UnsolvedMath contributors retain attribution for the source extract. Codex (GPT-6; exact serving revision not exposed) assisted with definitions, English summaries, scope notes and formatting. No new human mathematical authorship or independently verified solution is asserted. Catalogue identity checked against the complete supplied registry and current repository. Distinct upstream records, including repeated questions and broad or damaged extracts, are retained for complete coverage of the source CS category. A repeated record is not a new independent conjecture. Complete mathematical-equivalence and current-literature review has not been performed. Related canonical identities: 30001202. Preparation and final review follow contribution prompts C and D, with source gaps disclosed. Structural and rendering checks do not establish a new theorem or open status.
\label{end:statement}

\section{Literature and context check}
The primary report on printed page 673 contains the exact forward/backward recursions and a conjecture that they contain the finest information. The parent Codex GPT-6 supplied the concatenation proof route and the linear hyperbolic example; they have been checked and expanded here. The set-theoretic statement is settled below for arbitrary deterministic $f$, while the requested general conditions for smallness are addressed only by explicit sufficient regimes. This is not a general classification of nonlinear recovery. The original damaged introductory fragment and its link to 30001202 are retained. The argument is the same full auxiliary analysis as for that related record, not a new independent problem or a silent reconstruction of the fragment.
\section{Research attempt: Exact feasible-state projection and quantitative special regimes}

\paragraph{The exact information content of both recursions}
Let $\mathcal T$ be the set of all finite sequences $(x_0,\ldots,x_T)$ with $x_{t+1}=f(x_t)$ and $x_t\in P(t)$ for every $t$. No regularity or injectivity of $f$ is needed in this paragraph. Induction shows $Q^+(t)$ is precisely the set of endpoints of compatible prefixes: the assertion is true at zero, and taking $f$ of the preceding endpoint set followed by intersection with $P(t)$ is exactly the existential predecessor condition. Backward induction likewise shows $Q^-(t)$ is the set of starting points of compatible suffixes, because its inverse image uses all predecessors.

If $z\in Q^+(t)\cap Q^-(t)$, choose a prefix ending at $z$ and a suffix starting there. Their common state permits concatenation into a sequence in $\mathcal T$. Conversely any sequence in $\mathcal T$ supplies both witnesses. Therefore
\[
 Q(t)=\{x_t:(x_0,\ldots,x_T)\in\mathcal T\}.
 \tag{1}
\]
In particular $Q(t)$ is the smallest guaranteed state enclosure based solely on these exact observations and the known deterministic map: excluding any point of it would exclude one realizable full trajectory. This is a precise set-theoretic meaning of finest information, not a claim about efficient representation, probabilities, or approximate computations. Empty feasible sets mean inconsistent data; a realized unknown trajectory makes all of them nonempty.

\paragraph{Two sufficient diameter estimates}
Use Euclidean diameter. If $f$ is $L$-Lipschitz on the states being compared, with $0\leq L<1$, then for two feasible trajectories $\lVert x_t-y_t\rVert\leq L^t\lVert x_0-y_0\rVert$ by repeated application. Equation (1) thus gives
\[
 \operatorname{diam}Q(t)\leq L^t\operatorname{diam}P(0).
\]
Separately suppose $\lVert f(x)-f(y)\rVert\geq\lambda\lVert x-y\rVert$ for all relevant pairs, with $\lambda>1$. Applying this inequality to the suffix of two feasible trajectories gives
\[
 \operatorname{diam}Q(t)\leq\lambda^{-(T-t)}\operatorname{diam}P(T).
\]
This expansion assumption implies injectivity on the compared set; local derivative expansion alone is not a substitute for it when folding is possible. The source only requires the observed trajectory to stay in its compact observation region, not that every state be mapped back into that region.

\paragraph{An exact mixed stable/unstable example}
For $f(x,y)=(a x,b y)$ with $0<a<1<b$ and all observations $P(t)=[-1,1]^2$, a state $(x,y)$ at time $t$ has a compatible prefix precisely when $|x|\leq a^t$ and $|y|\leq1$. It has a compatible suffix precisely when $|x|\leq1$ and $|y|\leq b^{-(T-t)}$. These statements follow directly by iterating $a$ and $b$ in the respective time directions. Consequently
\[
 Q(t)=[-a^t,a^t]\times[-b^{-(T-t)},b^{-(T-t)}],\qquad
 \operatorname{diam}Q(t)=2\sqrt{a^{2t}+b^{-2(T-t)}}.
\]
Both sides shrink exponentially for times separated linearly from the two ends. This realizes the source's hyperbolic intuition without asserting a general stable-manifold theorem. By contrast the identity map with a repeated nontrivial cell has $Q(t)=P(t)$, so arbitrarily many observations need not shrink the set at all.

\section{Verification and adversarial audit}
The concatenation proof does not invert $f$ or assume a unique predecessor. Both endpoints $t=0,T$, inconsistent observations, constant maps, and the identity map are covered. Diameter estimates compare two complete feasible trajectories, so no independently chosen time slices are mistaken for a trajectory. The hyperbolic map need not preserve the whole compact square; only the admissible finite trajectories remain in it, as allowed in the primary source.
\paragraph{Audit corrections.} Separated exact information optimality from smallness and computational tractability; retained both repeated IDs and their original source qualifications.
\paragraph{External computation.} None; symbolic arguments and boundary checks only.
\section{Scope and remaining gaps}
Equation (1) answers the finest-enclosure subquestion under an explicit exact-data interpretation. The smallness results are sufficient assumptions and an exact linear example, not necessary-and-sufficient conditions for general nonlinear maps, efficient set algorithms, or a resolution of the broad/damaged source agenda. The parent model's mathematical suggestions are credited; no novelty claim is made.
\paragraph{Integration review.} Codex (GPT-6; exact serving revision not exposed) reviewed the written arguments and expanded layout and source-reference wrappers. Repository metadata and the production web renderer were checked using Python 3.13 and JavaScript with MathJax 3.2.2. This model-assisted review is not independent human proof verification; no PDF compilation or formal proof checking is asserted.

\section{Completion Estimate}
\noindent\textbf{COMPLETION: 0\%}
The calculations below the preserved source statement investigate a concrete route,
but establish no substantiated new progress toward the entire original target.
Known special cases and failed sufficient conditions are not credited as a solution.
\section{References}
The original source references and catalogue attribution are preserved in the
Sources section above. The following URLs were actually consulted for this attempt;
retention of an additional original reference does not claim fresh consultation of it.
\begin{itemize}
\item \url{https://ems.press/content/serial-article-files/46211} (consulted 2026-10-11). 
\end{itemize}
\section{Model and contributor disclosure}
The parent Codex GPT-6 model has reviewed the written mathematical arguments at their declared scopes. This is model-assisted review, not independent human verification.
OpenAI gpt-6-astra, requested reasoning configuration ultra, generated this research
attempt on 2026-10-11. The exact serving revision is not exposed. Codex GPT-6
supplied the mathematical suggestions credited above and performed orchestration and repository-format review; its exact serving revision
is also not exposed. The mathematical route, boundary checks and final adversarial
self-audit followed contribution prompts A then B in
\texttt{docs/contribution\_prompts.md}, after reading \texttt{CONTRIBUTING.md}.
Python 3.14.2 assembled the standalone documents; any mathematical computation is
specified in Verification. No theorem prover was used. No human mathematical
authorship, substantial human proof correction, or independent human proof checking
is asserted. Model-assisted review, source retrieval, and structural checks do not
establish correctness or novelty. Full private reasoning transcripts are unavailable
as reproducibility artifacts; the complete arguments, inputs and outcomes actually
used are stated here or in the companion batch audit.
\section{Final status}
Unresolved as a new research contribution. This is a concrete route investigation,
not a claim of new complete resolution or a current-status certification. A known
answer reproduced for context does not become a new contribution. The catalogue's
source-status label is preserved; this attempt status does not assert that a
source-labelled solved question is open.
\end{document}
